\section{Related Work}
\label{sec:related}

We survey prior work along four \emph{coverage} axes (taxonomy, engine, joint throughput and tail
reporting, and vendor involvement) spanned by this experiment (Supplement Table~S51), and separately
audit which comparability controls each source reports. A structured scoping search, not a systematic
review, of the ACM Digital Library, IEEE Xplore, Scopus, and Google Scholar (2006--July~2026)
retained empirical comparisons of relational access layers or engines, with details archived in the
replication package; being scoping rather than exhaustive, we scope every gap claim to the sources
searched.

\subsection{Object-relational mapping and the access-layer spectrum}
Access layers bridge the object-relational impedance mismatch between the relational model and an
application's object graphs~\cite{ireland2009impedance}, spanning native drivers, query builders, and
ORMs following Data Mapper / Active Record patterns~\cite{fowler2002patterns}
(Section~\ref{sec:introduction}). Torres et al.~\cite{torres2017orm} survey two decades of ORM
patterns; their trade-off (productivity against a hard-to-predict runtime cost) motivates what
we measure while holding the store family (relational rather than non-relational) fixed. The treatment is the configured
access-layer implementation bundle, including its selected API, SQL, protocol, and mapping path;
the study does not isolate a library binary from those choices.

\subsection{ORM performance and data-access anti-patterns}
A parallel line studies ORM performance in reverse, repairing the inefficient SQL frameworks emit.
The costliest is the N+1 query problem (Section~\ref{sec:introduction}); it and related
redundant-query anti-patterns pervade real-world web
applications~\cite{yang2018dbapps}, their impact quantified in ORM-backed
systems~\cite{chen2016redundant,chen2014orm} and often \emph{removable} by query synthesis,
batching, or caching. Closest to our
design, Lorenz et al.~\cite{lorenz2017orm} show the best mapping strategy depends on database
technology, van Zyl et al.~\cite{vanzyl2009optimisations} that it is tuning-sensitive, and a recent
study adds an energy dimension~\cite{bonvoisin2024orm}. Overhead is thus configurable, not fixed, 
as our per-pattern results confirm, but this line never compares drivers, query builders, and ORMs
head-to-head under load, our aim.

Cross-implementation ORM equivalence checking has a direct precedent. CYNTHIA generates an abstract
query, translates it to equivalent queries for five ORMs, executes them over four relational DBMSs,
and differentially compares the results to find correctness bugs~\cite{sotiropoulos2021cynthia};
matching results alone do not establish correctness against an independent task specification.
CYNTHIA is a testing method, not a performance-benchmark admission protocol: it does not use
equivalence to decide which timing cells may enter a benchmark, define a premeasurement rule for
which implementation represents each library, or separate equal-demand capacity effects from
matched-utilization latency. Our claim is restricted to using specification evidence and
cross-implementation equivalence as benchmark admission, joined with predeclared treatment selection
and operating-point separation.

\subsection{Peer-reviewed access-layer comparisons and adjacent lines}
Peer-reviewed head-to-head access-layer measurement is scarce, and each study leaves part of the
space uncovered. Colley et al.~\cite{colley2018orm} illustrate Entity Framework-generated against
hand-written SQL rather than comparing layers; two recent studies compare Prisma with optimized raw
SQL, and Prisma with TypeORM and Sequelize, both only on
PostgreSQL~\cite{prisma_rawsql_2025,orm_compare_jcct_2025}; and a third times internal query stages
of three ORMs rather than client-observed requests~\cite{zhadko2025orm}. Pratama and
Raharja~\cite{pratama2023nodebench} do vary five access products, but framework, access path, and
environment change across cells too, so no access-layer estimand is isolated. A further Express
study optimizes one stack rather than comparing layers~\cite{imamsakti2025async}.

A second, disjoint line benchmarks the layers below and above the access layer without varying it.
Salunke and Ouda~\cite{salunke2024pgmysql} compare PostgreSQL and MySQL under standard OLTP
workloads, the only dual-engine precedent found, but benchmark the \emph{engines themselves}. On the
Node.js side, work establishes Node's web viability~\cite{chaniotis2015nodejs}, tracks back-end drift
under sustained load~\cite{kuffel2025nodejs}, and compares frameworks
broadly~\cite{carmo2024backend}, in each case holding the access layer fixed.

\subsection{Standard benchmarks and benchmarking methodology}
Standard benchmarks define how systems are measured: YCSB~\cite{cooper2010ycsb} for cloud-serving
stores and TPC-C/TPC-E for OLTP. These target the \emph{engine}, not the driver, query builder, or ORM
this study varies, but set the bar for controlled, repeatable load. A methodological literature
explains why single-metric reporting fails: Fruth et al.~\cite{fruth2022tail} and Raasveldt et
al.~\cite{raasveldt2018fair} catalog benchmarking pitfalls, Dean and Barroso~\cite{dean2013tail} show
the tail drives user-perceived latency under fan-out, and Li et al.~\cite{li2014tales} trace its
sources. Such measurement is often ad hoc~\cite{leitner2017exploratory} and hard to reproduce, motivating a harness built for relevance, repeatability, and fairness. Following that guidance, we report
throughput \emph{and} p99 for the same run: a pairing the general literature recommends, which we did
not find in the access-layer benchmarks surveyed here.

\subsection{Practitioner and vendor benchmarks}
A larger body comes from practitioners and vendors, each authored by a party with a stake in a
compared product: Drizzle's Northwind micro-suite and k6 load test~\cite{drizzle_benchmarks},
Prisma's median-latency comparison reporting neither throughput nor any
percentile~\cite{prisma_benchmarks}, and IMDBench over three CRUD operations~\cite{imdbench}.
Together these span more products than the peer-reviewed literature, but none combines
dual-backend-stack coverage, response-time p99, capacity characterization, and frozen
source-located treatment provenance.

\subsection{Positioning}
\paragraph{Audit corpus.} The nine audited sources are those retained by the same scoping search,
under stated criteria rather than convenience: empirical performance comparisons of relational
access layers, or direct PostgreSQL-versus-MySQL comparisons, bearing on the four positioning
properties. Peer-reviewed studies, vendor suites, and practitioner benchmarks were all eligible,
since the audit concerns what a benchmark reports, not who published it. Coding was performed once,
by the author, from a stable locator for every judgement; \texttt{not\_reported} records what the
inspected source does not state, never a finding that the authors omitted the practice. The corpus
is not exhaustive, so every gap statement below is scoped to it: we do not claim that no other
benchmark implements these controls. Eligibility and exclusion criteria, the codebook, the
screening record, and per-judgement evidence are in the online supplement and the replication
package.

In the sources searched, no access-layer benchmark reports all seven protocol dimensions,
unsurprisingly, since the seven are this paper's own construction: the audit measures distance from a
new standard, not failure to meet a known one. The source-located audit records partial controls
rather than treating non-reporting as proof of absence (Supplement Table~S37), and we found no study
combining broad product coverage, both engines, joint throughput and tail reporting, and non-vendor
authorship, each covering at most two or three (Supplement Table~S51): the strongest claim the
scoping search supports, not one of priority. 

But coverage is not the contribution: the protocol operationalizes established experimental principles for
this genre, not a new one. Kounev et al.'s systems-benchmarking
criteria~\cite{kounev2020benchmarking}, Hasselbring's software-engineering
standard~\cite{hasselbring2021benchmarking}, and Raasveldt and M\"uhleisen's fair-benchmarking
guidance~\cite{raasveldt2018fair} supply its principles at a deliberately general level. They leave an
access-layer study to operationalize \emph{what code represents a library}, how competing
implementations are checked for semantic equivalence, and how strategy-sensitive results are
contrasted. Our protocol supplies one concrete operationalization: an arbitrary but predeclared treatment policy; a common-SQL raw-path sensitivity contrast; and a per-cell gate requiring specification-conformant, mutually equivalent responses and verified write state \emph{before} a timing is admitted. Neither element is new in itself, correctness-before-timing having precedents in SPEC validation, TPC ACID, and SQLancer~\cite{rigger2020sqlancer}, and equivalence checking in CYNTHIA above; the contribution is their composition into a discipline for establishing and reporting comparability among competing implementations of one task, stated in Section~\ref{sec:introduction}. Its evidence is the in-experiment ablation and a single-rater retrospective audit of nine external benchmarks (Supplement Table~S37); this demonstrates utility and applicability, not independent inter-rater validation or invention of the underlying principles.

